% Curated from paper commit 7ee9157; original equations and protocol identities retained.
\section{Branch conflict, ancestry routing and coefficient learning}
\label{note:conflict_ancestry}

\subsection{Context-gated Fashion-MNIST protocol}

The context-gated task used Fashion-MNIST classes 0 and 6. Inputs were
average-pooled to $7\times7$, standardized using training data and augmented
with a constant coordinate. Each seed used
8,192 training and 2,000 test draws with replacement from the official binary
splits, balanced across context and label. Branch counts were $B=2,4,8$;
conflict doses were $0,0.25,0.5,4/7,0.6,2/3,0.75,1$. Thresholding one
underlying uniform mask made conflict nested across doses. Correct, shared,
cyclically deranged, BP and gated-point rules reused inputs and initialization.
Full-batch descent used 250 updates at step 0.04 and an initial weight
standard deviation (SD) of 0.01.
Three exploratory seeds (100--102) selected a stable setting; an additional
implementation-check seed was also excluded. Analysis seeds 11000--11019
were fixed before outcomes, with no outcome-dependent stopping.

The primary estimand was each seed's least-squares slope of
correct-minus-shared accuracy on conflict, separately for each $B$.
Prespecified criteria required positive slopes in at least 16/20 seeds,
at least ten percentage points of separation at full conflict, less than
two points at zero conflict, correct routing above derangement and exact
BP/gated-point agreement. All criteria passed. The maximum analytic-versus-
autodifferentiation discrepancy was $3.43\times10^{-7}$; equivalent-rule
accuracy differences were zero. Intervals used 20,000 paired-seed bootstrap
draws; tests were two-sided Wilcoxon tests. Mean-field assumptions and empirical
chance crossings are distinguished below.

The same twenty seeds were replayed for post hoc gradient diagnostics at epochs 0, 1, 10, 50, 100 and 250, comparing both each method's own state and a common exact-rule state. Official input-file checksums were verified; the input manifest and all 28,800 rows are retained in \texttt{source\_data/review\_branch\_trajectories/}. Endpoint drift was below $5\times10^{-9}$ in loss. At conflict 0.5, common-state shared-gradient cosine fell from 0.987 to $-0.519$ for four branches and from 0.952 to $-0.701$ for eight branches by epoch 250 (Supplementary Fig.~\ref{fig:si_branch_conflict_controls}). The initialization boundary below was retained without refitting; trained gradient direction and trained accuracy crossing are separate endpoints.

\subsection{Mean-field boundary for context-gated branch selection}

The Fashion-MNIST conflict experiment selects one branch $c$ uniformly from
$B$ branches. Only $\bm w_c^{\mathsf T}\bm x_c$ enters the logit, so the
exact derivative of every unselected branch is zero. Shared feedback omits
that selector and updates every active eligibility. For the centered
mean-field prediction, assume balanced class signs $s\in\{-1,1\}$,
class means $\E[\bm x\mid s]=s\bm v$, and zero logits, so the logistic
error is $\delta=-s/2$. An unselected input copies the selected image with
probability $1-\chi$ and is independently drawn from the opposite class
with probability $\chi$. Its mean class evidence relative to the selected
view is therefore $1-2\chi$. The resulting compatibility matrix is
\[
C_\chi=(1-2\chi)\bm 1\bm 1^{\mathsf T}+2\chi I.
\]
Its shared-mode eigenvalue is $B-2\chi(B-1)$, which vanishes at
$\chi_c=B/[2(B-1)]$. Under these population and initialization assumptions,
the expected shared update loses positive alignment with the exact update
at that boundary. Logistic derivatives and class-conditional moments need
not preserve this relation as weights change. The empirical trained
chance crossing is consequently a separate descriptive endpoint, not a
theorem about the entire optimization trajectory. The missing branch
selector may be supplied in feedback or in local eligibility; it is not
evidence that the unselected branches require opposite exact errors.

\subsection{Two-branch subtree-address experiment}

The minimal two-sibling experiment used ten paired seeds (4100--4109), following an excluded numerical-check seed 4099. The check verified nondegenerate shared capture, exact-route capture one and analytic exact/BP agreement. The subsequent eight-stream factorial varied route budget.

For seed-specific unit vectors $\bm u_A$ and $\bm u_B$, binary label sign
$y\in\{-1,1\}$ and context $c\in\{0,1\}$, the two input means were
\begin{align}
\mathbb E[\bm x_A\mid y,c]
&=y\,[\mathbb 1(c=0)-0.75\mathbb 1(c=1)]\bm u_A,\\
\mathbb E[\bm x_B\mid y,c]
&=y\,[\mathbb 1(c=1)-0.75\mathbb 1(c=0)]\bm u_B.
\label{eq:s_credit_reversal_task}
\end{align}
Each stream had 12 coordinates and independent Gaussian noise with standard
deviation 1.15. The 2,048-example train and test splits contained exactly 512
examples in each context--label cell. The somatic logit selected branch $A$
when $c=0$ and branch $B$ when $c=1$. Both inputs nevertheless remained
active, preventing the one-coordinate rule from acquiring branch specificity
through a zero eligibility on the distractor.

For somatic logit $z$ and logistic sigmoid $\sigma$, let
$\delta=\sigma(z)-\mathbb 1(y=1)$ and
$\bm r(c)=(\mathbb 1[c=0],\mathbb 1[c=1])$. Exact and correct-subtree
coefficients were $\delta\bm r(c)$; within-neuron derangement reversed the
two entries, giving $\delta\bm r(c)_{::-1}$; the shared neuronal and same-depth controls used
$\delta(1,1)$; and the random rank-two control used
$\delta\bm r(c)R_s^{\mathsf T}$ for a seed-fixed orthogonal rotation $R_s$.
The gated point emulation used one somatic coefficient together with the
context-gated features $(\mathbb 1[c=0]\bm x_A,
\mathbb 1[c=1]\bm x_B)$, making its alternative address resource explicit.
All conditions had 24 trainable weights and the same paired initialization and
batch order. The random dense control had four routing nonzeros; the sparse
two-route conditions had two.

Training used mini-batch stochastic gradient descent (SGD), learning rate 0.3,
batch size 64 and 60 epochs.
The switch endpoint applied another 12 epochs to 1,024 context-1 examples with
balanced labels and measured the change in context-0 accuracy. For gradient
geometry, exact and candidate per-example gradients were concatenated over
both branches. Scaled capture was squared cosine after optimal scalar
rescaling. The one-step candidate was normalized to the exact batch-gradient
norm and evaluated at step 0.25. The Source Data indexed in Table~\ref{tab:source_index} retain every
tested condition; no condition was removed on the basis of its outcome.

All controls, including poor deranged and variable random-map outcomes, are retained. Correct routing, exact transport, BP and gated-point emulation coincide because the exact path coefficient is the binary selector. Thus this experiment isolates the supplied address resource; the gated-point control reproduces it.

\subsection{Subtree-address bandwidth and architecture factorial}

To vary feedback bandwidth, the complete factorial extended the two-sibling
task to eight context-selected streams arranged in a balanced binary hierarchy. Context and binary
label were exactly balanced in train and test sets, every stream remained
active on every example, and the sign of the task signal reversed outside the
selected context. Thus a shared coefficient could not recover the required address simply
because off-route eligibilities were zero. The prespecified feedback budgets
were $K\in\{1,2,4,8\}$ nested routes. Twenty paired analysis seeds
(4200--4219) crossed five representations and 27 method--budget combinations,
giving 135 fits per seed and 2,700 fits in total.

The dendritic tree, explicitly gated point model, flat-compartment model and
grouped-point model were implementation-equivalent controls: each contained 64
trainable weights, 64 active input contacts and one output nonlinearity, and
each received the same routed coefficient field. The fifth representation was
a degree- and depth-matched rewired tree. Its resource counts were unchanged,
but its ancestry groups were permuted relative to the task hierarchy. Feedback
families comprised correct ancestry, fixed within-neuron derangement,
depth-interleaved groups, random sparse groups, a seed-fixed random rank-$K$
map, and a task-fitted rank-$K$ upper bound. Neuron-shared, exact-transport and
analytic-backpropagation references completed the design. Route fields were
unit-normalized per example. Data, initialization and every
stochastic control map were keyed to seed, method family and $K$, so matched
representations saw identical stochastic realizations.

Pre-analysis checks required balanced labels and contexts, nonzero activity in every stream, shared capture below 0.2, exact capture one, analytic exact/BP equality and no fallback. An initial 2,700-fit cohort used stochastic maps unpaired across representation labels and was excluded. Only the complete repeat with paired maps and cross-representation equivalence checks enters inference.

The primary bandwidth contrasts are retained in the Source Data indexed in Table~\ref{tab:source_index}.
Correct nested-subtree routing strongly exceeded route derangement at $K=2$ and $K=4$, but
its comparison with the best matched non-anatomical control was deliberately
more stringent: ancestry lost at $K=1$ and $K=2$, won by 0.0127 accuracy at
$K=4$ (95\% paired seed-bootstrap interval 0.00593--0.01946; 15/20 seeds;
paired Wilcoxon signed-rank Holm-adjusted $P=0.0101$ across the four budgets;
the Source Data retain the paired mean-sign-flip sensitivity test), and tied at full rank. Correct routing on the
task-matched tree exceeded the rewired tree at $K=2$ and $K=4$ in all 20
seeds. The four matched representations were exactly equal in every condition
(maximum held-out-accuracy difference zero). The experiment therefore supports
an intermediate-bandwidth advantage for a task-aligned nested address, while
showing that an explicitly routed point implementation can reproduce it.

\subsection{Trained partition-residual reconstruction}

Using the saved configurations and outputs, we quantified irreducible address
error in a post hoc diagnostic of the complete factorial. For example $b$, let
$c_b$ denote context, $\delta_b$ the somatic error, $e_{c_b,i}$ the indicator
selecting leaf $i$, and $\bm x_{bi}$ its input vector. We defined the exact
coefficient field $c_{bi}=\delta_b e_{c_b,i}$, coordinate weight
$w_{bi}=\lVert\bm x_{bi}\rVert_2^2$, and active route field $\phi_{bi}$.
Allowing the best single coefficient for that route on each example gives
\begin{equation}
 z_b^*=\frac{\sum_i w_{bi}\phi_{bi}c_{bi}}
 {\sum_iw_{bi}\phi_{bi}^2},\qquad
 \mathcal C_{\mathrm{part}}=1-
 \frac{\sum_{b,i}w_{bi}(c_{bi}-z_b^*\phi_{bi})^2}
 {\sum_{b,i}w_{bi}c_{bi}^2}.
\label{eq:s_trained_partition_capture}
\end{equation}
This is the normalized complement of the shared-coefficient residual in
Eq.~\ref{eq:s_partition_variance}. It grants an optimal route coefficient and
therefore isolates address/support error from the coefficient scale actually
used during training.

All 20 seeds and 135 conditions per seed were reconstructed at initialization
and after 80 training epochs. Reconstructed losses agreed with the recorded
values within $4.97\times10^{-9}$ and accuracies within $5.01\times10^{-11}$,
the precision of the saved decimal tables; the address-plus-coefficient Pythagorean
decomposition closed within $7.28\times10^{-12}$. On the dendritic-tree
representation, trained correct-ancestry capture was 0.162, 0.315, 0.547 and
1.000 at $K=1,2,4,8$. Fixed coordinate-to-arbor derangement had the common scalar value
at $K=1$ and zero capture at $K=2,4,8$, because its active support excluded the
target leaf (Supplementary Fig.~\ref{fig:si_ancestry_coefficients}A).

For each seed we averaged implementation-equivalent representations within
feedback family and budget, then correlated trained capture with held-out
accuracy across conditions. The mean within-seed Spearman correlation was
0.824 (95\% seed-bootstrap interval 0.810--0.838; 20/20 positive; two-sided
Wilcoxon $P=1.91\times10^{-6}$). Restricting the analysis to the six route
families and $K<8$ removed the exact full-rank boundary and gave 0.639
(0.610--0.665; 20/20 positive; $P=1.91\times10^{-6}$; Source Data, Table~\ref{tab:source_index}). This association is descriptive
and post hoc. In particular, ancestry, depth-interleaved and random sparse
partitions had similar address capture at matched $K$, so the residual does
not explain the small trained ancestry advantage at $K=4$.

\subsection{Learning coefficients from an explicit local context cue}
\label{note:local_context_encoder}

To test whether coefficients could be learned from a supplied cue, we used a
four-route dictionary on fresh versions of the balanced eight-context task.
We retained its 64 forward weights, paired initialization,
2,048-example train and test sets, and 80 full-batch task-training updates.
The dictionary was fixed before task training:
$\Phi_{bk}=\mathbf{1}[k=\lfloor b/2\rfloor]/\sqrt2$ for zero-indexed
blocks $b=0,\ldots,7$ and routes $k=0,\ldots,3$.
For trial $i$ with context $c_i$, the supplied local cue was
$\bm h_i=\bm e_{c_i}+\sigma_h\bm\epsilon_i$, with
$\bm\epsilon_i\sim\mathcal N(0,I_8)$. Here $\bm e_{c_i}$ is the one-hot
context vector and $\sigma_h$ is cue-noise SD. The cue is an additional
information resource; the experiment does not infer a biological signal.

Separate balanced calibration trials supplied the target route-activation
vector $\bm e_{\lfloor c_i/2\rfloor}$. A four-output softmax encoder with
an intercept learned this mapping by minibatch delta-rule updates. Weights
started at zero; calibration used 30 epochs, batch size 16, step 0.2 and
weight decay 0.001 on non-bias weights. Calibration sets of 16, 64 and 256
trials were nested within each cue-noise condition. Because epoch count was
fixed, increasing sample count also increased the number of encoder updates.
This axis therefore changes both calibration data and computation.
No task labels, exact task gradients or held-out task outcomes entered encoder
training. Its parameters were frozen before the classification weights were
trained.

At cue delays $d=0,1,4$ trials, the encoder received the cue array circularly
shifted by $d$. Predicted coefficients $\bm c_i$ produced the normalized
field $\Phi\bm c_i/\|\Phi\bm c_i\|$. Delay is a mismatch between the
current eligibility and an earlier independent context; it is not a calibrated
physiological time constant. Controls used exact one-hot route selection,
a uniform frozen profile, or a cyclic permutation of the learned coefficient
coordinates. Exact and frozen controls repeated across nuisance conditions
are algebraic repetitions rather than independent evidence. Task updates used
the original factor-eight context rescaling and step 0.3.

Twenty analysis seeds (52000--52019) followed the excluded implementation-check
seed 51999.
Crossing three calibration sizes, cue-noise SDs $0,0.5,1$, delays and four
rules produced 2,160 task fits and 10,800 checkpoint rows. At task epochs
0, 1, 10, 40 and 80, we recorded held-out accuracy/loss and delivered-versus-
exact gradient cosine at each method's own state. Held-out cue classification,
coefficient mean squared error (MSE) and route-norm errors supplied separate
encoder diagnostics.
The protocol, fixed before new outcomes but not publicly preregistered, is
\texttt{configs/review\_completion/coefficient\_encoder.json}.

The primary condition used 256 calibration trials, cue-noise SD 0.5 and no
delay. Learned coefficients improved accuracy over the frozen profile by
7.16 percentage points (95\% paired-seed bootstrap interval 6.26--8.07;
20/20 positive) and over mismatched coefficients by 7.10 points
(6.19--7.99; 20/20). They remained 54.47 points below exact route selection
(53.54--55.45; 20/20 lower). Two-sided Wilcoxon tests restored the exact
$1/2048$ accuracy-count grid before ranking and used Holm correction across
these three primary contrasts. Adjusted $P$ values were
$5.72\times10^{-6}$ versus the frozen rule and
$1.77\times10^{-4}$ versus each other control. Intervals used 20,000
paired-seed draws; other grid contrasts are descriptive sensitivities.
All fits were finite. Learning from the supplied cue improved performance
relative to fixed or misassigned coefficients, but left a large gap to the
oracle. The experiment does not establish autonomous discovery of credit
coordinates in a biological neuron. The complete outputs are in
\texttt{source\_data/review\_coefficient\_encoder/}.

An exploratory follow-up was specified after observing the soft-encoder
outcomes. The same trained encoder, original paired seeds and task data were
retained, but its maximum-probability route was delivered as a one-hot vector.
No encoder parameter was refitted and no hyperparameter was selected from
the task outcomes. This hard readout is a paired sensitivity, not a new
independent replication or a replacement for the primary soft rule
(Supplementary Fig.~\ref{fig:si_ancestry_coefficients}). At the primary
noise SD 0.5, mean accuracy was 25.825\% for soft coefficients, 64.031\%
for hard selection, 80.291\% for exact selection and 18.669\% for the
frozen profile. Hard-minus-oracle accuracy remained $-16.26$ percentage
points (95\% paired-seed interval $-17.83$ to $-14.83$), while hard-minus-
frozen was 45.36 points (43.69--46.92). At zero noise and 256 calibration
trials, hard selection equaled the oracle numerically, whereas the soft rule
reached 79.978\%. With only 16 noise-free calibration trials, the soft
encoder already identified every context's route correctly, yet its graded
output gave 19.392\% task accuracy; selecting its maximum-probability route
gave the oracle's 80.291\%. This sensitivity isolates leakage into incorrect
routes from failure to identify the correct route. All hard-readout intervals and
unadjusted tests are descriptive; these reused conditions do not define a
new primary family. Both learned readouts failed when cues were delayed across
independent contexts. Thus accurate route identification does not guarantee
that graded coefficients are suitable for the learning rule: readout and
timing are separate constraints. All exploratory fits and unchanged encoder
settings are in \texttt{source\_data/review\_coefficient\_hard\_readout/}.

\subsection{Route overlap and task interference}

We finally examined how overlap between addressed regions and leakage outside
an intended route determine interference with an already learned task.
Branch-specific plasticity during motor learning suggests that different
tasks can recruit different dendritic regions \citep{cichon2015branch}. Let
$\bm w_A$ be an optimum of a quadratic Task-A objective with Hessian $H_A$,
and let an approximate Task-B teaching field produce update
$-\eta M_B\bm g_B$, where $\eta$ is the step size, $\bm g_B$ is the Task-B
gradient and $M_B$ is its routing operator. Expansion about the Task-A optimum removes the first-order
term, leaving the exact loss increase
\begin{equation}
\mathcal L_A(\bm w_A-\eta M_B\bm g_B)-\mathcal L_A(\bm w_A)
=\frac{\eta^2}{2}\bm g_B^\mathsf{T}M_B^\mathsf{T}H_AM_B\bm g_B.
\label{eq:supp_interference}
\end{equation}
In the equal-route example, opposing updates on the shared fraction $\omega$
have magnitude two, whereas leakage onto the Task-A-only fraction has
magnitude $\lambda$. Dividing by route size gives
$\eta^2[4\omega+\lambda^2(1-\omega)]/2$. Numerical vector updates matched this
identity to $4.2\times10^{-17}$ absolute error
(Eq.~\ref{eq:supp_interference}). Separated routes therefore
protect a learned task only when both anatomical overlap and teaching-field
leakage are small. The published somatostatin (SST) perturbation is consistent with the
resulting prediction that loss of branch-selective inhibition increases
effective leakage and cross-task interference \citep{cichon2015branch}, but it
does not uniquely identify the credit-routing mechanism. This static identity
does not model dendritic calcium, spikes, eligibility timing or the SST intervention.
